\medskip

\textbf{Partie A}

\medskip

On considère la fonction $f$ définie sur l'intervalle $[0~;~+\infty[$ par 
\[f(x) = \sqrt{x + 1}.\]

On admet que cette fonction est dérivable sur ce même intervalle.

\medskip

\begin{enumerate}
\item Démontrer que la fonction $f$ est croissante sur l'intervalle $[0~;~+\infty[$.
\item Démontrer que pour tout nombre réel $x$ appartenant à l'intervalle $[0~;~+\infty[$ :

\[f(x) - x = \dfrac{-x^2 + x + 1}{\sqrt{x + 1} + x}.\]

\item En déduire que sur l'intervalle $[0~;~+\infty[$ l'équation $f(x) = x$ admet pour unique solution :

\[\ell = \dfrac{1 +\sqrt 5}{2}.\]

\end{enumerate}

\medskip

\textbf{Partie B}

\medskip

On considère la suite $\left(u_n\right)$ définie par $u_0 = 5$ et pour tout entier naturel $n$, par 

$u_{n+1} = f\left(u_n\right)$ où $f$ est la fonction étudiée dans la \textbf{partie A}.

On admet que la suite de terme général $u_n$ est bien définie pour tout entier naturel $n$.

\medskip

\begin{enumerate}
\item Démontrer par récurrence que pour tout entier naturel $n$, on a 

\[1 \leqslant u_{n+1} \leqslant u_n.\]

\item En déduire que la suite $\left(u_n\right)$ converge.
\item Démontrer que la suite $\left(u_n\right)$ converge vers $\ell = \dfrac{1 +\sqrt 5}{2}$.
\item On considère le script Python ci-dessous :
\begin{center}
\begin{tabular}{|>{\ttfamily}l >{\ttfamily}l|}\hline
1 &\textbf{from} math \textbf{import} *\\
2 &\textbf{def} seuil(n):\\
3 &\qquad u=5\\
4 &\qquad i=$0$\\
5 &\qquad $\ell$ =(1 + \textbf{sqrt}(5))/2\\
6 &\qquad \textbf{while abs}(u-$\ell$)>=10**(-n):\\
7 &\qquad \qquad u=\textbf{sqrt}(u+1)\\
8 &\qquad \qquad i=i+1\\
9 &\qquad \textbf{return}(i)\\ \hline
\end{tabular}
\end{center}

\emph{On rappelle que la commande} \textbf{abs(x)} \emph{renvoie la valeur absolue de $x$}.

	\begin{enumerate}
		\item Donner la valeur renvoyée par {\ttfamily seuil} (2).
		\item La valeur renvoyée par {\ttfamily seuil} (4) est 9.
		
Interpréter cette valeur dans le contexte de l'exercice.
	\end{enumerate}
\end{enumerate}
